\scanpage{39}
\margin{\uline{Cor.} à lem\textsuperscript{2}} Si donc $H^0(X_0,\underline O_{X_0})=k$, cette indétermination ne dépend \unsure{plus|que} de $k$ : elle est levée en exigeant que une \deleted{pt}\inserted{section} marquée $s$ de $X$ aille sur \deleted{\unsure{une section} la}\inserted{une} section \unsure{marquée|unique} \deleted{[\unsure{supposant} que $S_{n-1}$ soit \unsure{sur} $\unsure{e_{n-1}}$].} // D'où en l'occurrence la possibilité de \unsure{fixer|faire|fixe} l'extension, avec la condition que la section $e\times e$ vienne en $e$.

De plus, \deleted{on \unsure{prend} sur $G$}\inserted{\unsure{cons. alors|considérons} $eg$} et $ge$, elles transforment $e$ en $e$ donc $e$ est \deleted{\unsure{déjà} section unité} d'après ce qui précède $eg=g$, $ge=g$ : $e$ est section unité bilatère. Pour la $=$ raison, on voit que $gg'=g'g$ et $(gg')g''=g(g'g'')$, car on a deux morphismes $G\times G\rightrightarrows G$ (resp\textsuperscript{t.} $G\times G\times G\rightrightarrows G$) qui coïncident \unsure{en réd. mod,|en réduction mod,|sur le mod,} et transforment $e\times e$ (resp\textsuperscript{t.} $e\times e\times e$) en $e$.

\noindent\rule{1.5cm}{0.3pt}

\uline{Cor.1 à lemme 3} Possibilité \inserted{et unicité} de \deleted{faire} l'extension dans le cas où on a la structure abélienne donnée sur $G_0$, et la section unité $e$ marquée sur $G$ [il suffit d'exiger $e\times e\longrightarrow e$], lorsque l'anneau est un \unsure{an} \uline{artinien}

\uline{Cor 2} Dans le cas \uline{formel} sur un anneau local complet.
% archive sheet number: 38
